\chapter{Introducción}
\label{ch:intro}

\section{Problema}

El discurso de odio y la toxicidad en redes sociales constituyen un desafío creciente
para la moderación de contenidos, en especial en contextos dialectales poco
representados en los corpus de entrenamiento masivos. Este informe documenta el
desarrollo y la evaluación de un sistema de clasificación multiclase de toxicidad
aplicado a publicaciones de redes sociales de El Salvador (plataformas X/Twitter y
Facebook).

El problema se formula como clasificación ordinal de cuatro niveles de gravedad, donde
el modelo debe distinguir entre contenido no tóxico, lenguaje ofensivo, discurso de
odio y amenazas o violencia.

\section{Esquema de clases}

Las etiquetas son ordinales (0--3) y se resumen en la Tabla~\ref{tab:clases}.

\begin{table}[H]
\centering
\caption{Esquema de clases del corpus.}
\label{tab:clases}
\begin{tabular}{cl}
\toprule
Label & Clase \\
\midrule
0 & No Tóxico \\
1 & Lenguaje Ofensivo \\
2 & Discurso de Odio \\
3 & Amenazas/Violencia \\
\bottomrule
\end{tabular}
\end{table}

\section{Objetivos y métrica principal}

El objetivo principal del proyecto es entrenar un modelo \emph{mBERT}
(\texttt{bert-base-multilingual-cased})~\cite{devlin2019bert} mediante fine-tuning
sobre el corpus local y alcanzar un F1-macro de al menos \textbf{0{,}80} en el
conjunto de prueba, según la meta definida en la configuración del proyecto.

Además se buscan:
\begin{itemize}[leftmargin=*]
  \item Establecer baselines clásicos (TF-IDF + clasificadores lineales).
  \item Evaluar robustez por cortes (plataforma, jerga, sarcasmo, ironía).
  \item Medir transferencia zero-shot a corpora externos en español.
  \item Explorar interpretabilidad con LIME~\cite{ribeiro2016lime} y
        SHAP~\cite{lundberg2017shap}.
\end{itemize}

\section{Resumen ejecutivo de resultados}

Sobre el conjunto de prueba ($n = 460$), el modelo mBERT alcanza un \textbf{F1-macro de 0{,}7777} (intervalo bootstrap 95\,\%: $[0{,}7358,\ 0{,}8144]$), superando al mejor baseline TF-IDF (LinearSVC, F1-macro $0{,}7645$). El objetivo de $0{,}80$ no se alcanza, aunque el intervalo de confianza superior lo incluye.


En evaluación zero-shot sobre seis datasets externos, el F1-macro cae al rango
$0{,}14$--$0{,}31$, lo que indica transferencia limitada fuera del dominio local.
Los cortes más difíciles son ironía (F1 $\approx$ 0.50) y sarcasmo (F1 $\approx$ 0.62), mientras que el desempeño en la plataforma X es sólido (F1 $\approx$ 0.83).


Los capítulos siguientes siguen el orden de los notebooks reproducibles del
proyecto (01--07): análisis exploratorio, preprocesamiento, baselines, fine-tuning,
evaluación completa, benchmark externo e interpretabilidad.
